\chapter{食物网}\label{ch:g5:food-webs}

\Cref{ch:g3:food-chains}提醒过我们，它那些整齐的\omterm{def:g3:food-chains:chain}{食物链}
只是从网上扯下的单根线：鹰也捕食老鼠，
蜥蜴也吃苍蝇，草也养活兔子。
这一章我们不再扯线，而是看整张
网——即\omterm{def:g5:food-webs:web}{食物网}——以及维持其平衡的因素。

\section{从食物链到食物网}

\begin{definition}[食物网]\label{def:g5:food-webs:web}
某一\omterm{def:g2:where-animals-live:habitat}{栖息地}的\emph{食物网}\index{食物网}是
由其\emph{全部}\omterm{def:g3:food-chains:chain}{食物链}同时组成的网络：每个\omterm{def:g5:biodiversity:species}{物种}
只画一次，箭头（仍表示\emph{\omterm{def:g3:food-chains:chain}{被……吃}}，
见\cref{rem:g3:food-chains:arrow}）表示每一种取食关系，
大多数\omterm{def:g1:animals-around-us:animal}{动物}取食多个\omterm{def:g5:biodiversity:species}{物种}，
也被多个\omterm{def:g5:biodiversity:species}{物种}吃，因此\omterm{def:g3:food-chains:chain}{食物链}相互交叉，
图形便成了一张网。
\end{definition}

\begin{omfigure}
\begin{tikzpicture}[scale=0.98]
  % nodes
  \node[draw, thick, omMeth, fill=omMeth!8, rounded corners=3pt,
    font=\small] (grass) at (0,0) {草};
  \node[draw, thick, omMeth, fill=omMeth!8, rounded corners=3pt,
    font=\small] (clover) at (2.6,0) {三叶草};
  \node[draw, thick, omMeth, fill=omMeth!8, rounded corners=3pt,
    font=\small] (oak) at (5.2,0) {橡树叶};
  \node[draw, thick, omExo, fill=omExo!8, rounded corners=3pt,
    font=\small] (hopper) at (0,1.7) {蚱蜢};
  \node[draw, thick, omExo, fill=omExo!8, rounded corners=3pt,
    font=\small] (rabbit) at (2.6,1.7) {兔子};
  \node[draw, thick, omExo, fill=omExo!8, rounded corners=3pt,
    font=\small] (caterpillar) at (5.2,1.7) {毛毛虫};
  \node[draw, thick, omProp, fill=omProp!8, rounded corners=3pt,
    font=\small] (lizard) at (0,3.4) {蜥蜴};
  \node[draw, thick, omProp, fill=omProp!8, rounded corners=3pt,
    font=\small] (tit) at (5.2,3.4) {蓝山雀};
  \node[draw, thick, omThm, fill=omThm!8, rounded corners=3pt,
    font=\small] (fox) at (2.6,4.4) {狐狸};
  \node[draw, thick, omThm, fill=omThm!8, rounded corners=3pt,
    font=\small] (buzzard) at (0.2,5.2) {鹰};
  \node[draw, thick, omThm, fill=omThm!8, rounded corners=3pt,
    font=\small] (hawk) at (5.0,5.2) {雀鹰};
  % arrows
  \draw[->, thick, omDef] (grass) -- (hopper);
  \draw[->, thick, omDef] (grass) -- (rabbit);
  \draw[->, thick, omDef] (clover) -- (rabbit);
  \draw[->, thick, omDef] (oak) -- (caterpillar);
  \draw[->, thick, omDef] (hopper) -- (lizard);
  \draw[->, thick, omDef] (caterpillar) -- (tit);
  \draw[->, thick, omDef] (rabbit) -- (fox);
  \draw[->, thick, omDef] (lizard) -- (buzzard);
  \draw[->, thick, omDef] (hopper) -- (tit);
  \draw[->, thick, omDef] (tit) -- (hawk);
  \draw[->, thick, omDef] (rabbit) -- (buzzard);
\end{tikzpicture}

{\small 草地边缘的一张小型\omterm{def:g5:food-webs:web}{食物网}。
\cref{ch:g3:food-chains}中的每一条单线\omterm{def:g3:food-chains:chain}{食物链}都包含在内，
且彼此交织。}
\end{omfigure}

\begin{example}[读懂食物网]\label{ex:g5:food-webs:reading}
看图可以读出两点。顺着兔子的箭头看：它
既吃草\emph{和}三叶草，又被狐狸\emph{和}
鹰吃掉——共四种关系，一条链只显示两种。
再数一数鹰的食物箭头：蜥蜴和兔子，
去掉一个，另一个仍在。这张网恰恰记录了
\cref{prop:g5:biodiversity:web}所称道的那些备选途径。
\end{example}

\section{网中的角色}

\begin{definition}[生产者与消费者]\label{def:g5:food-webs:roles}
网中的\omterm{def:g5:biodiversity:species}{物种}可分为以下几类角色：
\begin{itemize}
  \item \emph{生产者}\index{生产者}：\omterm{def:g1:plants-around-us:plant}{植物}，
        它们自己制造物质
        （见\cref{prop:g4:what-plants-need:make}）——每条
        箭头路径都始于生产者
        （见\cref{prop:g3:food-chains:start}）；
  \item \emph{消费者}\index{消费者}：\omterm{def:g1:animals-around-us:animal}{动物}，
        以其他\omterm{def:g1:living-or-not:living}{生物}为食——\omterm{def:g2:what-animals-eat:kinds}{植食性动物}直接取食生产者，
        \omterm{def:g2:what-animals-eat:kinds}{肉食性动物}取食其他消费者，
        \omterm{def:g2:what-animals-eat:kinds}{杂食性动物}两者兼有（见\cref{def:g2:what-animals-eat:kinds}）；
  \item 还有在幕后默默工作的
        \emph{分解者}\index{分解者}——正是
        \cref{ex:g2:caring-for-nature:compost}中的那些腐烂者，它们
        以动\omterm{def:g1:plants-around-us:plant}{植物}遗体为食，把其中的物质归还给
        土壤。以后的年级会专门为它们设一章。
\end{itemize}
\end{definition}

\begin{example}[图中生物的分类]\label{ex:g5:food-webs:sorting}
在该小型\omterm{def:g5:food-webs:web}{食物网}中：\omterm{def:g5:food-webs:roles}{生产者}——草、三叶草、橡树；
植食性\omterm{def:g5:food-webs:roles}{消费者}——蚱蜢、兔子、\omterm{ex:g2:animals-grow-up:butterfly}{毛毛虫}；
肉食性\omterm{def:g5:food-webs:roles}{消费者}——蜥蜴、蓝山雀（主要）、狐狸
（实为\omterm{def:g2:what-animals-eat:kinds}{杂食性动物}）、鹰、雀鹰。三层结构：
绿色层、植食层、肉食层——每一上层
都站在其下层之上。
\end{example}

\begin{proposition}[数量的金字塔形状]\label{prop:g5:food-webs:pyramid}
逐层统计网中居民的数量，得到一个
稳定的形态：\omterm{def:g5:food-webs:roles}{生产者}\emph{很多}，
\omterm{def:g2:what-animals-eat:kinds}{植食性动物}\emph{较少}，\omterm{def:g2:what-animals-eat:kinds}{肉食性动物}\emph{很少}——一片草地上
数以百万计的草株养活数千只蚱蜢，蚱蜢又养活
一小群蜥蜴，蜥蜴又养活一对鹰。每一层
都依赖其下一层生存，数量不能超过其食物。这
就是为什么大型猎食者总是十分稀少。
\end{proposition}
\admitted

\section{平衡，以及打破平衡的因素}

\begin{proposition}[食物网的平衡]\label{prop:g5:food-webs:balance}
在无外力干扰时，\omterm{def:g5:food-webs:web}{食物网}能让各类\omterm{def:g1:living-or-not:living}{生物}的数量大致保持
\emph{平衡}\index{平衡（食物网）}：若兔子增多，
它们的猎食者吃得更好，也随之增多，
多余的数量便会消退；若兔子减少，饥饿的猎食者
便转而捕食其他猎物（即备用的箭头），兔子数量
随之回升。这些调节作用
（见\cref{prop:g3:food-chains:depend}）沿着许多条箭头同时发生，
共同使整个网保持稳定。
\end{proposition}
\admitted

\begin{example}[经典的错误]\label{ex:g5:food-webs:mistake}
农民们曾发起运动，要消灭猛禽——
“偷鸡贼”。运动成功之处，老鼠和田鼠
失去主要天敌后大量繁殖成鼠害，
吃掉的粮食远超过鹰偶尔抓走的鸡的价值。
园丁讲的蜗牛故事
（见\cref{exo:g3:food-chains:9}），放大到整个乡村的尺度来看：移除
\omterm{def:g5:food-webs:web}{食物网}的顶层，其下一层
便会猛增。
\end{example}

\begin{method}[预测扰动的影响]\label{met:g5:food-webs:predict}
要预测移除（或增加）某\omterm{def:g5:biodiversity:species}{物种}的后果，可按以下步骤进行：
\begin{enumerate}
  \item 在\omterm{def:g5:food-webs:web}{食物网}中找到该\omterm{def:g5:biodiversity:species}{物种}，列出其\emph{所有}箭头，
        包括进入与离开的箭头；
  \item 其食物因被捕食减少，数量往往增加；其捕食者
        因食物减少，数量往往下降，或转而捕食其他猎物；
  \item 沿\omterm{def:g5:food-webs:web}{食物网}再追踪每一变化一步——
        第二步的影响往往出人意料；
  \item 检查哪些邻近\omterm{def:g5:biodiversity:species}{物种}有备用箭头（它们能吸收冲击），
        以及哪些\omterm{def:g5:biodiversity:species}{物种}只依赖
        此\omterm{def:g5:biodiversity:species}{物种}（它们处于危险之中）。
\end{enumerate}
\end{method}

\section{练习}

\begin{exercise}[$\star$]\label{exo:g5:food-webs:1}
什么是\omterm{def:g5:food-webs:web}{食物网}？为什么各条\omterm{def:g3:food-chains:chain}{食物链}会交织在一起？
\end{exercise}

\begin{exercise}[$\star$]\label{exo:g5:food-webs:2}
说出\cref{def:g5:food-webs:roles}中的三种角色，
各举一例，取自图中（\omterm{def:g5:food-webs:roles}{分解者}除外）。
\end{exercise}

\begin{exercise}[$\star$]\label{exo:g5:food-webs:3}
根据图示：列出蓝山雀所吃的一切食物，
以及吃蓝山雀的所有\omterm{def:g1:living-or-not:living}{生物}。
\end{exercise}

\begin{exercise}[$\star$]\label{exo:g5:food-webs:4}
从图中找出两条完整的\omterm{def:g3:food-chains:chain}{食物链}，它们
都包含蚱蜢。
\end{exercise}

\begin{exercise}[$\star$]\label{exo:g5:food-webs:5}
逐层说明数量的形状。为什么任何一层的数量
都不能超过其下一层？
\end{exercise}

\begin{exercise}[$\star$]\label{exo:g5:food-webs:6}
用一句话回答：为什么鹰很少，而蚱蜢
多得数不清？
\end{exercise}

\begin{exercise}[$\star$]\label{exo:g5:food-webs:7}
兔子迎来丰年，大量繁殖。依据
\cref{prop:g5:food-webs:balance}，接下来会发生什么——以及
什么使多余的数量消退？
\end{exercise}

\begin{exercise}[$\star$]\label{exo:g5:food-webs:8}
动手试一试：画出一个你熟悉的\omterm{def:g2:where-animals-live:habitat}{栖息地}的\omterm{def:g5:food-webs:web}{食物网}——如花园、
公园、池塘——至少包含六个\omterm{def:g5:biodiversity:species}{物种}，且每条箭头
都能说出理由。用绿色标出\omterm{def:g5:food-webs:roles}{生产者}。
\end{exercise}

\begin{exercise}[$\star\star$]\label{exo:g5:food-webs:9}
对移除图中蜥蜴的情形，运用\cref{met:g5:food-webs:predict}
进行分析：先看第一步的影响，再追踪一步。
\end{exercise}

\begin{exercise}[$\star\star$]\label{exo:g5:food-webs:10}
用一连串的预测重述\cref{ex:g5:food-webs:mistake}：
若移除鹰，接下来会怎样，再之后又会怎样？
发起运动者没有预见到的是哪一步？
\end{exercise}

\begin{exercise}[$\star\star\star$]\label{exo:g5:food-webs:11}
比较图中\omterm{def:g5:food-webs:web}{食物网}里移除橡树（\omterm{def:g5:food-webs:roles}{生产者}）
与移除雀鹰（顶级猎食者）的情况。哪些
邻近\omterm{def:g5:biodiversity:species}{物种}最先感受到各自损失的影响，每一扰动
沿哪个方向传播，哪一种损失在这里
更容易被\omterm{def:g5:food-webs:web}{食物网}吸收？请运用
\cref{ex:g3:food-chains:dry,ex:g3:food-chains:hunter}。
\end{exercise}
